% Демонстрация 4: «Сечения поверхностей 2-го порядка» (Лекция 11)
\section*{Демонстрация 4. «Сечения поверхностей 2-го порядка»\\ \normalsize\textmd{Файл \texttt{demos/04-quadric-sections.html} \quad(к Лекции~11)}}

\subsection*{Какие факты лекции иллюстрирует}

Демонстрация делает наглядным центральный инструмент Лекции~11 ---
\textbf{метод параллельных сечений}, с помощью которого исследуется
форма невырожденных поверхностей второго порядка. В интерактиве
доступны эллипсоид, одно- и двуполостный гиперболоиды и эллиптический
параболоид; ползунки полуосей $a$, $b$, $c$ задают каноническое
уравнение, а ползунок $h$ перемещает секущую плоскость $z = h$, под
3D-видом немедленно перерисовывается линия сечения.

\begin{itemize}
  \item \textbf{Канонические уравнения.} При движении ползунков панель
    формул показывает уравнение поверхности с подставленными
    значениями: для эллипсоида
    $\dfrac{x^2}{a^2} + \dfrac{y^2}{b^2} + \dfrac{z^2}{c^2} = 1$, для
    однополостного гиперболоида
    $\dfrac{x^2}{a^2} + \dfrac{y^2}{b^2} - \dfrac{z^2}{c^2} = 1$, для
    двуполостного --- та же левая часть, равная $-1$, для параболоида
    $z = \dfrac{x^2}{a^2} + \dfrac{y^2}{b^2}$. Студент видит, как знаки
    при квадратах координат управляют типом поверхности.

  \item \textbf{Метод сечений в формулах.} Подставляя $z = h$, для
    эллипсоида получаем эллипс
    \[
      \frac{x^2}{a^2} + \frac{y^2}{b^2} = 1 - \frac{h^2}{c^2},
    \]
    который сжимается в точку при $|h| \to c$. Для однополостного
    гиперболоида сечение $z = h$ --- эллипс
    $\dfrac{x^2}{a^2} + \dfrac{y^2}{b^2} = 1 + \dfrac{h^2}{c^2} \ge 1$,
    наименьший («горловой») при $h = 0$; вертикальные сечения дают
    гиперболы. Для двуполостного гиперболоида правая часть
    $\dfrac{h^2}{c^2} - 1$ неотрицательна лишь при $|h| \ge c$, и
    демонстрация явно показывает «пустое» сечение в полосе $|z| < c$.
    Эти выкладки полностью совпадают с разбором определений из
    Лекции~11 (эллипсоид, оба гиперболоида).

  \item \textbf{Влияние полуосей на форму.} Раздельное изменение $a$,
    $b$, $c$ показывает, как вытягивается и сплющивается поверхность и
    как меняются полуоси эллипса сечения; равенство $a = b = c$ для
    эллипсоида даёт сферу.

  \item \textbf{Прямолинейные образующие.} Для однополостного
    гиперболоида включается отдельный флажок «Показать прямолинейные
    образующие»: на изогнутой поверхности появляются два семейства
    прямых. Это визуализация того, что однополостный гиперболоид и
    гиперболический параболоид --- \textbf{дважды линейчатые}
    поверхности (через каждую точку проходят две прямые образующие),
    тогда как эллипсоид, двуполостный гиперболоид и эллиптический
    параболоид образующих не имеют.
\end{itemize}

\begin{remarkIT}{Где это работает за пределами лекции}{demo04_apps}
  Сюжет демонстрации напрямую выходит на инженерные и
  IT-приложения:
  \begin{itemize}
    \item \textbf{Гиперболоидные конструкции.} Шуховские башни и
      градирни ТЭЦ собираются из прямых стальных балок, образующих
      жёсткую изогнутую оболочку, --- это и есть линейчатость
      однополостного гиперболоида, которую студент видит на экране.
    \item \textbf{Антенны и оптика.} Параболоид вращения фокусирует
      параллельный пучок в одну точку --- спутниковые «тарелки» и
      зеркала телескопов.
    \item \textbf{CAD и 3D-моделирование.} Линейчатые и седловые куски
      --- базовые «патчи» при построении поверхностей сплайнами; форма
      набирается из прямых элементов.
    \item \textbf{Метод поверхностей уровня (level-set).} Сечение
      $z = h$ --- это линия уровня функции двух переменных; тот же
      приём лежит в основе визуализации функций потерь и контурных
      карт в Data Science.
  \end{itemize}
\end{remarkIT}

\subsection*{Примерный сценарий использования на занятии}

\begin{enumerate}
  \item \textbf{Когда включить.} Демонстрацию запускают сразу после
    определения эллипсоида и формулировки метода сечений (определение
    \texttt{ellipsoid\_def} Лекции~11), как только на доске выписано
    уравнение сечения $\dfrac{x^2}{a^2} + \dfrac{y^2}{b^2} =
    1 - \dfrac{h^2}{c^2}$.

  \item \textbf{Эллипсоид и движение плоскости.} Выберите «Эллипсоид»,
    медленно ведите ползунок $h$ от $0$ к $c$. Вопрос аудитории:
    \emph{«Какой формы сечение при $z = 0$ и при $z$, близком к $c$?»}
    Студенты должны предсказать эллипс наибольшего размера в центре и
    его стягивание в точку у полюса; экран подтверждает ответ и
    одновременно показывает обновляющуюся правую часть формулы.

  \item \textbf{Роль полуосей.} Меняйте $a$ и $b$ по отдельности, затем
    сделайте $a = b = c$. Спросите, какая поверхность получилась
    (сфера) и почему все сечения стали окружностями.

  \item \textbf{Однополостный гиперболоид.} Переключитесь на
    однополостный гиперболоид и снова прогоните ползунок $h$. Обратите
    внимание: сечение всегда остаётся эллипсом и никогда не исчезает,
    наименьший --- «горловой» при $h = 0$. Свяжите это с неравенством
    $1 + \dfrac{h^2}{c^2} \ge 1$.

  \item \textbf{Двуполостный гиперболоид и пустое сечение.} Перейдите к
    двуполостному гиперболоиду и медленно увеличивайте $|h|$ от нуля.
    Вопрос: \emph{«Почему в некотором интервале высот сечения нет
    вообще?»} Демонстрация показывает пустое сечение при $|h| < c$ ---
    это полоса без точек поверхности, разделяющая две полости.

  \item \textbf{Момент удивления: прямые на кривой поверхности.}
    Вернитесь к однополостному гиперболоиду и включите флажок
    прямолинейных образующих. Дайте паузу: на явно изогнутой
    поверхности проступают \emph{прямые линии}, причём два семейства.
    Это и есть линейчатость, которую трудно вообразить без 3D-вида;
    свяжите с теоремой о наличии образующих и с факторизацией
    $\dfrac{x^2}{a^2} - \dfrac{z^2}{c^2} = 1 - \dfrac{y^2}{b^2}$ из
    Лекции~11. Здесь же уместно напомнить про Шуховские башни.

  \item \textbf{Параболоид как контрпример.} Выберите эллиптический
    параболоид: сечения $z = h$ --- эллипсы лишь при $h > 0$, при
    $h = 0$ вырождаются в вершину, при $h < 0$ исчезают. Подчеркните,
    что у параболоида образующих нет --- поверхность целиком лежит в
    полупространстве $z \ge 0$.

  \item \textbf{Связь с методом исследования поверхностей.}
    Подытожьте: вращая 3D-вид и двигая плоскость, мы фактически
    выполнили тот же алгоритм, что и на доске, --- по серии плоских
    сечений восстановили форму пространственной поверхности. Именно
    так метод параллельных сечений работает при ручном исследовании
    любой квадрики по её каноническому уравнению.
\end{enumerate}
